\documentclass{article}
\usepackage[paper=letterpaper,margin=2cm]{geometry}
\usepackage[utf8]{inputenc}
\usepackage[russian]{babel}
\usepackage[]{graphicx}
\usepackage[usenames]{color}
\usepackage{colortbl}
\usepackage{geometry}
\usepackage{xcolor}
\usepackage{hyperref}
\usepackage{../../lib/latex/listings-rust}
\usepackage{fontspec}
\setmonofont{JetBrains Mono}[Contextuals=Alternate,Ligatures = TeX,]
\usepackage{listings}
\usepackage{keycommand}
\usepackage{caption}

\setmainfont[
  Ligatures=TeX,
  Extension=.otf,
  BoldFont=cmunbx,
  ItalicFont=cmunti,
  BoldItalicFont=cmunbi,
]{cmunrm}
\setsansfont[
  Ligatures=TeX,
  Extension=.otf,
  BoldFont=cmunsx,
  ItalicFont=cmunsi,
]{cmunss}

\geometry{
  a4paper,
  top=25mm,
  right=30mm,
  bottom=25mm,
  left=30mm
}

\hypersetup{
  colorlinks=true,
  linkcolor=blue!50!red,
  urlcolor=blue!70!black
}

\captionsetup[lstlisting]{
  font={tt},
}

% based on Atom One Light
\lstset{
  language=Java,
  frame=single,
  basicstyle=\ttfamily\color[HTML]{383a42},
  columns=fullflexible,
  breaklines=true,
  numbers=left,
  frame=tab,
  postbreak=\mbox{\textcolor{red}{$\hookrightarrow$}\space},
  extendedchars=false,
  showspaces=false,
  showstringspaces=false,
  identifierstyle=\ttfamily\color[HTML]{4078f2},
  commentstyle=\color[HTML]{a0a1a7},
  stringstyle=\color[HTML]{50a14f},
  keywordstyle=\color[HTML]{a626a4},
  numberstyle=\ttfamily\color[HTML]{2c91af},
  rulecolor=\color[HTML]{383a42}
}

\lstdefinelanguage{XML}
{
  morestring=[b]",
  morestring=[s]{>}{<},
  morecomment=[s]{<?}{?>},
}

\newcommand{\code}[1]{
  \lstset{title=#1}
  \lstinputlisting{#1}
}
\newkeycommand{\itmo}[variant=aboba, labn=aboba, discipline=aboba, group=aboba, student=aboba,teacher=aboba, year=2022]{
  \begin{titlepage}
    \begin{center}
      \section*{
        Федеральное государственное автономное образовательное учреждение\\ высшего образования\\
        «Национальный исследовательский университет ИТМО»\\
        Факультет Программной Инженерии и Компьютерной Техники \\
       }
      \includegraphics[scale=0.2]{../../lib/img/itmo.png}
    \end{center}

    \vspace{4cm}

    \begin{center}
      \large \textbf{Вариант \textnumero \commandkey{variant}}\\
      \textbf{Лабораторная работа \textnumero \commandkey{labn}}\\
      по дисциплине\\
      \textbf{\commandkey{discipline}}
    \end{center}

    \vspace*{\fill}

    \begin{flushright}
      Выполнил Студент группы \commandkey{group}\\
      \textbf{\commandkey{student}}\\
      Преподаватель: \\
      \textbf{\commandkey{teacher}}\\
    \end{flushright}

    \vspace{1cm}

    \begin{center}
      г. Санкт-Петербург\\
      \commandkey{year}г.
    \end{center}

    \thispagestyle{empty}
  \end{titlepage}
}


\begin{document}

\itmo[
      variant=18,
      labn=4,
      discipline=Информатика,
      group=P3115,
      student=Владимир Мацюк,
      teacher=Малышева Татьяна Алексеевна
]
\lstset{language=Python}

\tableofcontents

\section{Задание}
\begin{enumerate}
      \subsection{Вариант}
      \item Определить номер варианта как остаток деления на 36 порядкового
            номера в списке группы в ISU. В случае, если в данный день недели
            нет занятий, то увеличить номер варианта на восемь.
            $$
                  \begin{tabular}{|c|c|c|c|c|c|c|c|c|}
                        \hline
                        18 & JSON & XML & Четверг \\
                        \hline
                  \end{tabular}
            $$
      \item  Изучить форму Бэкуса-Наура.
      \item  Изучить особенности языков разметки/форматов JSON, YAML, XML.
      \item  Понять устройство страницы с расписанием для своей группы: \url{http://itmo.ru/ru/schedule/0/P3110/schedule.htm}
      \item 
            \subsection{Исходный файл}
      \item  Исходя из структуры расписания конкретного дня, сформировать
            файл с расписанием в формате, указанном в задании в качестве
            исходного. При этом необходимо, чтобы в выбранном дне было не
            менее двух занятий (можно использовать своё персональное). В
            случае, если в данный день недели нет таких занятий, то увеличить
            номер варианта ещё на восемь.
            \lstinputlisting{input.json}
            
            \subsection{Обязательное задание}
      \item  Обязательное задание (позволяет набрать до 65 процентов от
            максимального числа баллов БаРС за данную лабораторную):
            написать программу на языке Python 3.x, которая бы осуществляла
            парсинг и конвертацию исходного файла в новый.
      \item  Нельзя использовать готовые библиотеки, в том числе регулярные
            выражения в Python и библиотеки для загрузки XML-файлов.
            \lstset{language=Python}
            \code{task1.py}
            \lstset{language=XML}
            \code{out1.xml}
            
            \subsection{Дополнительное задание \textnumero 1}
      \item  Дополнительное задание \textnumero 1 (позволяет набрать +10
            процентов от максимального числа баллов БаРС за данную
            лабораторную).
            \begin{enumerate}
                  \item Найти готовые библиотеки, осуществляющие аналогичный
                        парсинг и конвертацию файлов.
                  \item Переписать исходный код, применив найденные
                        библиотеки. Регулярные выражения также нельзя
                        использовать.
                  \item Сравнить полученные результаты и объяснить их
                        сходство/различие.
            \end{enumerate}
            
            \lstset{language=Python}
            \code{task2.py}
            \lstset{language=XML}
            \code{out2.xml}
            
            \subsection{Дополнительное задание \textnumero 2}
      \item Дополнительное задание \textnumero 2 (позволяет набрать +10
            процентов от максимального числа баллов БаРС за данную
            лабораторную).
            \begin{enumerate}
                  \item Переписать исходный код, добавив в него использование
                        2
                        регулярных выражений.
                  \item Сравнить полученные результаты и объяснить их
                        сходство/различие.
            \end{enumerate}
            
            
            \lstset{language=Python}
            \code{task3.py}
            \lstset{language=XML}
            \code{out3.xml}
            
            \subsection{Дополнительное задание \textnumero 3}
      \item Дополнительное задание \textnumero 3 (позволяет набрать +10
            процентов от максимального числа баллов БаРС за данную
            лабораторную).
            \begin{enumerate}
                  \item Используя свою исходную программу из обязательного
                        задания, программу из дополнительного задания \textnumero 1 и
                        программу из дополнительного задания \textnumero 2, сравнить
                        стократное время выполнения парсинга + конвертации в
                        цикле.
                  \item Проанализировать полученные результаты и объяснить их
                        сходство/различие.
            \end{enumerate}
            
            \lstset{language=Python}
            \code{task4.py}
            \lstset{language=Python}
            \code{out4.txt}
            
            \subsection{Дополнительное задание \textnumero 4}
      \item  Дополнительное задание \textnumero 4 (позволяет набрать +5
            процентов от максимального числа баллов БаРС за данную
            лабораторную).
            \begin{enumerate}
                  \item  Переписать исходную программу, чтобы она осуществляла
                        парсинг и конвертацию исходного файла в любой другой
                        формат (кроме JSON, YAML, XML, HTML): PROTOBUF,
                        TSV, CSV, WML и т.п.
                  \item  Проанализировать полученные результаты, объяснить
                        особенности использования формата.
            \end{enumerate}
            
            \lstset{language=Python}
            \code{task5.py}
            \lstset{language=Python}
            \code{out5.toml}
            
            \subsection{Итог}
      \item  Проверить, что все пункты задания выполнены и выполнены верно.
      \item  Написать отчёт о проделанной работе.
      \item  Подготовиться к устным вопросам на защите
\end{enumerate}

\section{Вывод}
sdfsdf

\end{document}
